\paragraph{Local existence and Taylor solutions.}
The existence question is answered for each invariant graph, rather than
inferred from a plotted numerical trajectory. For a regular third-order row
write its corrected pullback as
\[
 L(z,F,P,C)H+N(z,F,P,C)=0,\qquad
 P=F',\quad C=F'',\quad H=F''',\quad R=-N/L.
\]
At an initial jet $(z_0,F_0,P_0,C_0)$ satisfying the recorded positive-chart
conditions and $L_0\ne0$, the right-hand side of
\[
 (F,P,C)'=(P,C,R(z,F,P,C))
\]
is real analytic in a neighborhood. This remains true for fixed noninteger
$m$ on a strictly positive velocity branch. The analytic local ODE theorem
\cite[Section~4.1]{teschl2012}
therefore gives a unique analytic profile germ for the specified three
initial values. The verified similarity pullback then gives a local invariant
solution of the printed scalar master equation. Continuity preserves the
strict chart inequalities on a sufficiently small neighborhood. This
argument supplies neither a numerical convergence radius nor a global
solution, and does not establish the full bed-field model.

For $\xi=z-z_0$ put $F=\sum_{j\geq0}a_j\xi^j$, where
$a_0=F_0$, $a_1=P_0$ and $a_2=C_0/2$. If
$F_{n+2}=\sum_{j=0}^{n+2}a_j\xi^j$, coefficient matching gives
\begin{equation}
 a_{n+3}=
 \frac{[\xi^n]R(z_0+\xi,F_{n+2},F_{n+2}',F_{n+2}'')}
 {(n+3)(n+2)(n+1)},\qquad n\geq0.
 \label{eq:sve-series-recurrence}
\end{equation}
The recurrence constructs these coefficients successively. At order zero,
$F'''(z_0)=R_0$ gives $a_3=R_0/6$. Differentiating the normal form once
along its solution gives the next coefficient explicitly:
\[
 a_4=\frac1{24}\left(R_z+P_0R_F+C_0R_P+R_0R_C\right)_0.
\]
All derivatives on the right are evaluated at the same initial jet. Higher
coefficients follow by the same coefficient-matching rule; they are not
fitted to numerical values. Substitution of the resulting polynomial into
the corrected undivided ODE provides a separate algebraic check. For degree $N\geq3$, it checks the coefficients of
$\xi^0,\ldots,\xi^{N-3}$ exactly. The remaining ODE residual is
$O(\xi^{N-2})$; the local truncation errors in $F,F',F''$ are respectively
$O(\xi^{N+1})$, $O(\xi^N)$ and $O(\xi^{N-1})$. A finite polynomial with
this residual property is not normally an exact solution.

For common scaling, set $(m,g,\kappa)=(3,1,1/30)$ and use initial data
\[
 (z_0,F_0,P_0,C_0)=(1,2,1/2,1/10).
\]
Direct substitution gives
$L_0=79/20$, the remaining residual term is $N_0=11/100$, and $F'''(1)=-11/395$, hence
\[
 F(1+\xi)=2+\frac{\xi}{2}+\frac{\xi^2}{20}
             -\frac{11\xi^3}{2370}+O(\xi^4).
\]
This determines a local similarity solution $\phi(t,s)=tF(s/t)$, including
the initial point $t=s=1$. For row 3.8 with initial jet $(0,1,1/10)$ at
$t=0$, the series terminates: $F(t)=t+t^2/20$ is exact.

The following cubic data make the five regular solution families explicit
at $m=3$, $g=1$ and $\kappa=1/30$. Each row starts
$F(z_0+\xi)=F_0+P_0\xi+C_0\xi^2/2+a_3\xi^3+O(\xi^4)$;
the recurrence determines the subsequent coefficients. The nonzero values
of $L_0$ and positive chart margins justify the regular local construction.
\begin{center}\small
\begin{tabular}{@{}lcccc@{}}
\toprule
Case & $z_0$ & $(F_0,P_0,C_0)$ & $L_0$ & $a_3$\\
\midrule
3.4 & $1$ & $(2,1/2,1/10)$ & $79/20$ & $-11/2370$\\
3.5 & $1$ & $(2,1/2,1/10)$ & $-22/5$ & $83/1100$\\
3.6 & $1$ & $(2,1/2,1/10)$ & $-71/20$ & $169/21300$\\
3.8 & $0$ & $(0,1,1/10)$ & $1$ & $0$\\
3.9 & $1$ & $(2,1/2,1/10)$ & $-79/20$ & $173/23700$\\
\bottomrule
\end{tabular}
\end{center}

The other rows need different conclusions. Row 3.7 admits
$\phi=\alpha t+F(s)$ with arbitrary sufficiently smooth $F$, $\alpha>0$
and $F'>0$; its identity pullback cannot determine a unique Taylor
recurrence. Row 3.1 is stationary and lies outside the positive-velocity
chart. Row 3.2 violates $\phi_s>0$, and row 3.3 has no invariant graph.
A zero leading coefficient means the regular existence test is inconclusive,
not that no solution exists: the separately verified exact cubic power
profile lies on such a common-scaling normal-form singularity. Numerical
comparisons of Taylor truncations with the established IVPs are local
diagnostics, separate from the analytic-germ existence statement.

The degree-six campaign constructs all $21$ regular IVP series and checks
all four residual coefficients, from $\xi^0$ through $\xi^3$, exactly.
Comparisons use four saved RKF45/DOP853 grid nodes per IVP, without
interpolation, and retain separate errors for $F,F',F''$. For the cubic
common-scaling example, the profile differences from DOP853 are
$3.28\times10^{-9}$ at $\xi=0.2$ and $3.24\times10^{-13}$ at $\xi=0.05$.
At $\xi=0.025$ the difference is below the operational resolution threshold;
it is not another measured convergence order. The numerical comparisons
neither bound the series remainder nor locate its convergence radius.
